\documentclass[conference]{IEEEtran}
\IEEEoverridecommandlockouts
\usepackage{cite}
\usepackage{amsmath,amssymb,amsfonts}
\usepackage{graphicx}
\usepackage{textcomp}
\usepackage{xcolor}
\usepackage{booktabs}
\usepackage{multirow}
\usepackage{array}
\usepackage{url}
\usepackage{makecell}
\def\BibTeX{{\rm B\kern-.05em{\sc i\kern-.025em b}\kern-.08em
    T\kern-.1667em\lower.7ex\hbox{E}\kern-.125emX}}

\begin{document}

\title{Does More Context Help? Evaluating Knowledge-Augmented Prompt Strategies
for Arabic Post-OCR Correction with a Lightweight LLM}

\author{\IEEEauthorblockN{Mohamed Sabry}
\IEEEauthorblockA{\textit{Department of Electronics and Electrical Communications} \\
\textit{Faculty of Engineering, Cairo University}\\
Giza, Egypt \\
mohamedsabry109@gmail.com}
\and
\IEEEauthorblockN{Mohsen Rashwan}
\IEEEauthorblockA{\textit{Department of Electronics and Electrical Communications} \\
\textit{Faculty of Engineering, Cairo University \& RDI}\\
Cairo, Egypt \\
mrashwan@rdi-eg.ai}
}

\maketitle

\begin{abstract}
Arabic optical character recognition (OCR) systems produce systematic
character-level errors---dot confusions, hamza substitutions, taa marbuta
swaps---that degrade downstream text processing. Post-OCR correction using
large language models (LLMs) is a natural remedy, but it is unclear how much
additional context actually helps, and whether findings generalize across
font families, script types, and model scales.
This paper reports a three-experiment controlled study using
Qwen3-4B-Instruct-2507 as the primary correction model. Experiment~1
evaluates eight prompt-engineering strategies across nine datasets (eight
typewritten fonts from PATS-A01 and KHATT handwritten). Experiment~2
extends the evaluation to thirteen datasets, adding KHATT-Paragraph
(handwritten paragraphs), Yarmouk (contemporary printed documents),
Muharaf (handwritten manuscripts), and Historical (archival pages),
testing eight prompt configurations. Experiment~3 benchmarks four LLMs
across two OCR sources and two held-out benchmarks (RDI-Test-Lines,
Kitab). All metrics are computed on normal samples only (OCR/GT length
ratio $\leq 5$) with diacritics stripped before CER/WER computation.

In Experiment~1, we evaluate eight prompt-engineering strategies.
The central finding is that zero-shot LLM correction \emph{harms} the
PATS-A01 macro-average CER (5.57\%~$\to$~5.84\%, $+$4.8\%), degrading 7 of
8 fonts. Only domain-matched BM25 RAG (Phase~8) beats the OCR baseline
(5.45\%, $-$0.12~pp). On KHATT (34.24\% CER), all strategies improve by
0.4--1.1~pp; the best is Error-Signature RAG (33.13\%, $-$1.11~pp).

In Experiment~2, conservative zero-shot prompting (T3) achieves the best
PATS result: 5.27\% CER ($-$5.4\% relative), outperforming RAG without
retrieval overhead. Base+RAG (T2) is the best overall configuration (PATS
5.44\%, KHATT 33.82\%) and is used in Experiment~3.

In Experiment~3, four models are compared with T2 across multiple domains.
Gemma-3-12B achieves 5.33\% PATS CER; Qwen3-4B is best on KHATT (33.82\%)
and full-page datasets (OCR 86.50\%~$\to$~corrected 68.18\%).
LLM correction generalises to unseen benchmarks: RDI-Test-Lines $-$10.7\%,
Kitab Benchmark $-$17$\to$$-$24\% CER reduction.
These findings establish that over-correction is universal at low OCR quality,
that conservative prompting is more effective than knowledge augmentation for
clean typewritten text, and that RAG-anchored correction generalises to
diverse unseen domains.
\end{abstract}

\begin{IEEEkeywords}
Arabic OCR, post-OCR correction, large language models, prompt engineering,
retrieval-augmented generation, BM25, Qwen3, PATS-A01, KHATT, KHATT-Paragraph,
Yarmouk, Muharaf, historical manuscripts, multi-domain evaluation
\end{IEEEkeywords}

% -------------------------------------------------------------------------
\section{Introduction}
% -------------------------------------------------------------------------

Arabic optical character recognition occupies a uniquely difficult position
in the document analysis landscape. The script is inherently cursive, with
letter forms that vary by position within the word (isolated, initial, medial,
final); dots serve as the sole phonemic and semantic distinguisher for entire
character groups (\textit{ba}/\textit{ta}/\textit{tha}/\textit{ya}/\textit{nun}
share an identical skeletal form); and the optional nature of diacritics
introduces a pervasive ambiguity that complicates both recognition and
evaluation. These properties conspire to produce systematic, structured errors
in open-source OCR engines such as Qaari~\cite{qaari}---errors that propagate
unchecked into downstream pipelines for information retrieval, machine
translation, and text-to-speech synthesis.

Post-OCR correction---treating recognition output as noisy text and applying a
language model as a denoiser---has attracted renewed attention since
instruction-tuned large language models (LLMs) became widely available. The
paradigm is appealing: it requires no task-specific training data, operates in a
single forward pass, and leverages the model's implicit prior over valid
linguistic sequences. Yet two questions remain inadequately addressed in the
literature. First, \emph{does injecting structured knowledge about the OCR
engine's failure modes into the prompt yield measurable gains over zero-shot
correction?} Second, and more fundamentally, \emph{does correction effectiveness
generalise across font families and script types, or is it an artefact of the
specific evaluation conditions?}

To answer these questions, we design a controlled multi-dataset experiment that
starts from a zero-shot baseline and progressively augments the prompt with
orthogonal knowledge sources (a failure-filtered confusion matrix,
self-reflective diagnostics, morphological post-processing, and retrieval),
each evaluated in isolation against the zero-shot baseline. The central finding
is that correction effectiveness is governed by a \emph{font-dependent
over-correction threshold}: on moderate-to-high-error fonts (10--25\% CER)
zero-shot correction yields large reductions, but on the lowest-error font and
on segmentation-dominated handwritten text it is net harmful or marginal.

The contributions of this work are:
\begin{itemize}
\item A three-experiment systematic study across thirteen datasets spanning
      CER 2.12--129\%, revealing font- and domain-dependent behaviour that
      single-dataset studies structurally cannot detect.
\item A formal characterisation of the \emph{over-correction threshold}: the
      baseline CER below which LLM correction is net harmful, estimated at
      $\approx$6--7\% for Qwen3-4B on Qaari output.
\item A novel self-reflective prompting strategy that injects the model's own
      training-split error statistics and word-level failure pairs into the
      system prompt, giving the strongest isolated PATS gain ($-$0.7~pp CER)
      and uniquely improving the lowest-error font.
\item A domain-matched BM25 RAG strategy that retrieves similar correction
      pairs as per-sample few-shot context; domain alignment avoids the
      catastrophic degradation of external-corpus RAG and is the only
      Experiment~1 strategy to beat the OCR baseline on PATS-A01.
\end{itemize}

% -------------------------------------------------------------------------
\section{Related Work}
% -------------------------------------------------------------------------

Arabic OCR has been studied for over three decades~\cite{arabic_ocr_survey},
and its core difficulties are structural: the cursive baseline couples
segmentation with recognition, diacritics are contextually optional yet
semantically significant, and several letter groups (ba/ta/tha/ya/nun; fa/qaf;
dal/dhal; ra/zay) share identical skeletal forms differing only in dots. These
dot-group confusions are the dominant error type in modern engines and the
primary target of the confusion-matrix augmentation evaluated here. Qaari, the
open-source engine used throughout, is a vision--language model (VLM)
fine-tuned from Qwen2-VL-2B~\cite{qaari}; it exhibits systematic, font-dependent
dot-confusion patterns that are quantifiable.

Post-OCR correction has evolved through three paradigms. Rule-based and
statistical methods used lexicon lookup, $n$-gram models, and noisy-channel
formulations~\cite{afli2016}; they are effective for frequent patterns but
require manual curation and struggle with unseen errors. Neural
sequence-to-sequence approaches treat correction as translation over parallel
(OCR, GT) corpora~\cite{dong2018,nguyen2021}, but need substantial training data
and are tightly coupled to a specific engine's error distribution. LLM-based
methods---the paradigm explored here---correct OCR output via natural-language
prompts without task-specific training and can outperform specialised models on
English~\cite{llm_ocr_correction}, yet Arabic evaluations remain limited to
single-dataset BERT-class studies~\cite{arabic_bert_ocr}. No prior work
systematically evaluates multiple prompt augmentations across multiple Arabic
fonts under a controlled isolation design.

Knowledge injection into prompts yields consistent NLP gains: few-shot
exemplars activate task priors~\cite{brown2020}, chain-of-thought decomposes
reasoning~\cite{wei2022}, and retrieval-augmented generation grounds outputs in
external evidence~\cite{lewis2020}. Most relevant, Reflexion~\cite{shinn2023}
feeds an agent its own prior failures as verbal reinforcement; our Phase~4
self-reflective strategy adapts this to a single-turn correction task, giving
the model a statistical summary of its training-split errors and concrete
word-level failure pairs---to our knowledge the first application of
self-diagnostic prompting to post-OCR correction in any language. We further use
CAMeL Tools~\cite{obeid2020,habash2010} for morphological validation (Phase~5)
and error characterisation (Phase~1), and evaluate on the complementary
PATS-A01~\cite{pats} (eight fonts over identical text, enabling controlled
cross-font comparison) and KHATT~\cite{khatt} (real handwritten,
segmentation-dominated, far higher error rates) benchmarks.

% -------------------------------------------------------------------------
\section{Datasets and Experimental Setup}
% -------------------------------------------------------------------------

\subsection{Datasets}
The three experiments use complementary Arabic OCR datasets spanning
line-level and page-level documents across typewritten, handwritten,
and historical domains. \textbf{PATS-A01}~\cite{pats} is a typewritten dataset
of eight font families; we use an aligned 85/15 train/validation split
(seed=42), yielding 415 validation samples per font from a shared 2{,}766-line
pool, so that font is the sole cross-font variable. Training splits feed
self-reflective analysis (Phase~4); validation splits feed all reported metrics.
\textbf{KHATT}~\cite{khatt} is a handwritten dataset (34.24\% baseline CER),
a substantially harder challenge than any typewritten font. Four
\textbf{extended full-page datasets} (Experiment~2) test page-level
generalisation: \textbf{KHATT-Paragraph} (449 handwritten paragraph images),
\textbf{Yarmouk} (1{,}663 contemporary printed encyclopedic pages),
\textbf{Muharaf} (116 handwritten documents), and \textbf{Historical} (40
degraded archival manuscript pages). Table~\ref{tab:datasets} summarizes all
evaluation datasets.

\begin{table}[t]
\caption{Evaluation datasets. Exp1 = Experiment~1 (9 datasets).
Exp2 adds four full-page datasets (total 13). OCR CER = Qaari baseline
(normal samples, diacritics stripped). $\dagger$ After runaway correction.}
\label{tab:datasets}
\begin{center}
\small
\begin{tabular}{llrr}
\toprule
\textbf{Dataset} & \textbf{Type} & \textbf{Samples} & \textbf{OCR CER} \\
\midrule
\multicolumn{4}{l}{\textit{Experiment 1 datasets (line-level)}} \\
PATS-Akhbar       & Typewritten & 415 & 2.12\% \\
PATS-Simplified   & Typewritten & 415 & 2.94\% \\
PATS-Traditional  & Typewritten & 415 & 3.63\% \\
PATS-Arial        & Typewritten & 415 & 3.65\% \\
PATS-Naskh        & Typewritten & 415 & 4.10\% \\
PATS-Tahoma       & Typewritten & 415 & 5.82\% \\
PATS-Thuluth      & Typewritten & 415 & 10.21\% \\
PATS-Andalus      & Typewritten & 415 & 12.12\% \\
KHATT-val         & Handwritten & 186 & 34.24\% \\
\midrule
\multicolumn{4}{l}{\textit{Experiment 2 additional datasets (page-level)}} \\
KHATT-Paragraph   & HW paragraph  & 449  & 61.68\% \\
Yarmouk-testing   & Print (encycl.) & 1{,}663 & 49.85\% \\
Muharaf-val       & HW document   & 116  & 129.39\%$\dagger$ \\
Historical        & HW manuscript & 40   & 105.10\%$\dagger$ \\
\bottomrule
\end{tabular}
\end{center}
\end{table}

Diacritics are stripped before CER and WER evaluation to avoid penalizing the
model for diacritic insertion or deletion, which is not the target of this
correction task. Both diacritics-stripped and raw metrics are computed; we
report the stripped variant throughout.\footnote{Results with diacritics
retained show the same relative patterns.}

\subsection{OCR Engine and Correction Model}
We use Qaari as the OCR engine; its output serves as both the Phase~1 baseline
and the input to all LLM correction phases.
Qwen3-4B-Instruct-2507~\cite{qwen3}, a 4-billion-parameter instruction-tuned
model with native Arabic support, is the primary correction LLM. Inference ran
on a GPU-accelerated environment (Kaggle T4) with the \texttt{enable\_thinking}
flag set to \texttt{False} to suppress the chain-of-thought scratchpad, which
adds token cost without measurable benefit here.

\subsection{Evaluation Metrics}
Character Error Rate (CER) is the primary metric:
\begin{equation}
\text{CER} = \frac{S + D + I}{N},
\label{eq:cer}
\end{equation}
where $S$, $D$, $I$ are character-level substitutions, deletions, and
insertions, and $N$ is the reference character count. Word Error Rate (WER) is
reported alongside CER. Lower values indicate better performance. Each phase
follows a three-stage pipeline---\textit{export} (local), \textit{infer}
(cloud GPU with HuggingFace checkpointing), and \textit{analyze} (local against
ground truth)---decoupling correction from evaluation.

% -------------------------------------------------------------------------
\section{Experimental Methodology}
% -------------------------------------------------------------------------

The study is organized into the phases of Table~\ref{tab:methodology}. Phase~1
characterizes the OCR errors and builds the per-dataset confusion matrix and
morphological error taxonomy; Phase~2 is the zero-shot LLM correction baseline.
Phases~3, 4, 5, 8, and 9 each introduce exactly one knowledge source under an
\emph{isolation design}---changing one variable relative to Phase~2, so any
phase is beneficial only if it beats Phase~2's CER. Phase~6 tests combinations
of the isolated strategies with ablation. All correction prompts use a
conservative, XML-structured system template (``return as-is if correct; fix
only clear visual OCR errors; do not add diacritics''). The RAG phases (8, 9)
retrieve the top-$k=5$ most similar (OCR, GT) pairs from the Phase~2
\emph{training-split} corrections---a same-engine, same-font, same-domain
corpus---and inject them as a per-sample few-shot block; Phase~8 retrieves by
BM25~\cite{robertson2009bm25} over character $n$-grams ($n\in\{2,3,4\}$),
Phase~9 by error-signature similarity (CAMeL morphological profile $+$ confusion
profile). Preliminary external-corpus retrieval (OpenITI~\cite{romanov2017openiti})
degraded severely from domain mismatch, motivating the same-domain corpus.

\begin{table}[t]
\caption{Experimental phases. All correction phases use the conservative
zero-shot prompt as their base. Phase~5 results are excluded from comparisons
(implementation bugs; corrected re-implementation evaluated in Experiment~2).}
\label{tab:methodology}
\begin{center}
\footnotesize
\setlength{\tabcolsep}{4pt}%
\begin{tabular}{p{0.04\columnwidth}p{0.22\columnwidth}p{0.26\columnwidth}p{0.26\columnwidth}}
\toprule
\textbf{Ph.} & \textbf{Method} & \textbf{Knowledge source} & \textbf{Novel element} \\
\midrule
1 & OCR baseline + taxonomy & Qaari output, CAMeL & Confusion matrix, valid-but-wrong rate \\
2 & Zero-shot correction & none (conservative prompt) & Comparison baseline \\
3 & OCR-aware prompting & confusion matrix (P1) & Failure-filtered confusions \\
4 & Self-reflective & P2 training-split errors & Self-diagnostic stats + word pairs \\
5 & CAMeL post-proc. & morphological analyzer & \emph{excluded---implementation bugs} \\
6 & Combinations & P3+P4 (+P5) & Ablation of components \\
8 & RAG (BM25) & training corrections & Char $n$-gram retrieval \\
9 & Error-sig.\ RAG & CAMeL + confusion + training & Error-structural retrieval \\
\bottomrule
\end{tabular}
\end{center}
\end{table}

% -------------------------------------------------------------------------
\section{Results}
% -------------------------------------------------------------------------

\subsection{OCR Baseline and Zero-Shot Correction}

Table~\ref{tab:phase12} reports the Qaari OCR baseline and the Phase~2 zero-shot
correction for all nine Experiment~1 datasets. The OCR error rate varies
$5.7\times$ across PATS fonts (2.12--12.12\% CER, macro-average 5.57\%), so
typeface choice has a larger effect than any post-processing we test; KHATT
(34.24\%) is $6.1\times$ harder. Morphological analysis shows 71\% of OCR errors
are valid-but-wrong (real but contextually incorrect Arabic words) and 29\% are
non-words, bounding what a morphological revert can catch. Zero-shot correction
is net harmful on PATS (macro-average 5.57\%~$\to$~5.84\%, $+$4.8\%), degrading
seven of eight fonts; only Traditional benefits ($-$0.67~pp), and since
Traditional (3.63\%) is helped while Arial (3.65\%) is harmed, a simple CER
threshold cannot explain the pattern---error type and systematicity matter.
On KHATT, zero-shot gives a consistent but modest $-$0.99~pp, as segmentation
failures resist text-only correction.

\begin{table}[t]
\caption{OCR baseline (Phase~1) and zero-shot LLM correction (Phase~2),
diacritics stripped. $\Delta$CER is P2$-$OCR in percentage points; negative
indicates improvement. Fonts ordered by increasing OCR CER.}
\label{tab:phase12}
\begin{center}
\small
\begin{tabular}{lrrr}
\toprule
\textbf{Dataset} & \textbf{OCR CER} & \textbf{P2 CER} & \textbf{$\Delta$CER} \\
\midrule
PATS-Akhbar       &  2.12 &  2.71 & +0.59 \\
PATS-Simplified   &  2.94 &  2.93 & -0.01 \\
PATS-Traditional  &  3.63 &  2.96 & -0.67 \\
PATS-Arial        &  3.65 &  4.50 & +0.85 \\
PATS-Naskh        &  4.10 &  4.27 & +0.17 \\
PATS-Tahoma       &  5.82 &  6.13 & +0.31 \\
PATS-Thuluth      & 10.21 & 10.92 & +0.71 \\
PATS-Andalus      & 12.12 & 12.27 & +0.15 \\
PATS-A01 Macro-Avg&  5.57 &  5.84 & +0.27 \\
\midrule
KHATT-val         & 34.24 & 33.25 & -0.99 \\
\bottomrule
\end{tabular}
\end{center}
\end{table}

\subsection{Phases 3--9: Isolated Knowledge Augmentation}

Table~\ref{tab:isolated} compares all isolated strategies against the Phase~2
baseline, macro-averaging across PATS-A01 fonts with KHATT reported separately.
Phase~3 (confusion matrix) and Phase~4 (self-reflective) give small PATS gains
($-$0.10 and $-$0.25~pp vs.\ P2); Phase~4 is the strongest isolated PATS
strategy, reaching near-parity with the OCR baseline (5.59\% vs.\ 5.57\%), and
is uniquely able to improve the lowest-error font by explicitly warning the
model against over-correction. Phase~8 (domain-matched BM25 RAG) achieves the
best PATS-A01 result (5.45\%, $-$0.39~pp vs.\ P2, $-$0.12~pp vs.\ OCR)---the
only strategy to beat the OCR baseline---while Phase~9 (Error-Signature RAG)
achieves the best KHATT result (33.13\%, $-$1.11~pp). Phase~5 (CAMeL) is
excluded: post-hoc analysis found the morphological analyser was silently
disabled and the revert keyed by token rather than position, so its values
(21.89\% PATS, 41.62\% KHATT) are invalid.

\begin{table}[t]
\caption{Isolated augmentation strategies vs.\ Phase~2 zero-shot baseline.
PATS-A01 CER is the unweighted mean across all eight fonts; KHATT separate.
Diacritics stripped. Best LLM result per column in \textbf{bold}.
$\dagger$ Phase~5 excluded (implementation bugs).}
\label{tab:isolated}
\begin{center}
\small
\begin{tabular}{lcccc}
\toprule
& \multicolumn{2}{c}{\textbf{PATS-A01 (avg)}} & \multicolumn{2}{c}{\textbf{KHATT}} \\
\cmidrule(lr){2-3} \cmidrule(lr){4-5}
\textbf{Phase} & \textbf{CER} & \textbf{$\Delta$P2} & \textbf{CER} & \textbf{$\Delta$P2} \\
\midrule
1 (OCR baseline)         &  5.57 & ---    & 34.24 & ---   \\
2 (Zero-shot)            &  5.84 & +0.27  & 33.25 & -0.99 \\
\midrule
3 (OCR-Aware)            &  5.74 & -0.10  & 33.26 & -0.98 \\
4 (Self-Reflective)      &  5.59 & -0.25  & 33.24 & -1.00 \\
5 (CAMeL)\textsuperscript{$\dagger$} & 21.89 & n/a & 41.62 & n/a  \\
6 (conf\_self)           &  5.76 & -0.08  & 33.27 & -0.97 \\
\textbf{8 (RAG, BM25)}  &  \textbf{5.45} & \textbf{-0.39} & 33.83 & -0.41 \\
9 (Error-Sig RAG)        &  5.47 & -0.37  & \textbf{33.13} & \textbf{-1.11} \\
\bottomrule
\end{tabular}
\end{center}
\end{table}

\subsection{Phase 6: Combinations and Ablation}

Table~\ref{tab:combos} reports strategy combinations. RAG strategies dominate:
Phase~8 is best on PATS (5.45\%, the only strategy beating the OCR baseline) and
Phase~9 best on KHATT (33.13\%). The \texttt{conf\_self} combination (P3+P4) is
weaker on PATS (5.76\%) than P4 alone (5.59\%), indicating that adding
confusion-matrix context to an already self-aware model slightly increases
intervention on clean fonts. The Phase~8 ablation against zero-shot shows RAG
eliminating harm on the highest-error fonts (Andalus $-$1.86~pp, Thuluth
$-$1.24~pp) while only marginally worsening Akhbar---the domain-matched corpus
suppresses false corrections where they hurt most.

\begin{table}[t]
\caption{Phase 6: Strategy combinations vs.\ Phase~2 zero-shot baseline.
conf=confusion-matrix injection, self=self-reflective. PATS-A01 CER is the
unweighted macro-average. Diacritics stripped. Best per column in \textbf{bold}.}
\label{tab:combos}
\begin{center}
\small
\begin{tabular}{lcccc}
\toprule
& \multicolumn{2}{c}{\textbf{PATS-A01 (avg)}} & \multicolumn{2}{c}{\textbf{KHATT}} \\
\cmidrule(lr){2-3} \cmidrule(lr){4-5}
\textbf{Combination} & \textbf{CER} & \textbf{$\Delta$P2} & \textbf{CER} & \textbf{$\Delta$P2} \\
\midrule
P2 (zero-shot)           & 5.84 & ---   & 33.25 & ---   \\
conf\_only (P3)          & 5.74 & -0.10 & 33.26 & -0.01 \\
self\_only (P4)          & 5.59 & -0.25 & 33.24 & -0.01 \\
conf\_self (P3+P4)       & 5.76 & -0.08 & 33.27 & +0.02 \\
\textbf{P8 (RAG, BM25)} & \textbf{5.45} & \textbf{-0.39} & 33.83 & +0.58 \\
P9 (Error-Sig RAG)       & 5.47 & -0.37 & \textbf{33.13} & \textbf{-0.12} \\
\bottomrule
\end{tabular}
\end{center}
\end{table}

\subsection{Per-Font Analysis}

The best per-font strategy (Table~\ref{tab:perfont}) confirms that no single
approach dominates: \texttt{conf\_self} wins on the lowest-error fonts,
Error-Signature RAG on the mid-range, and BM25 RAG on the highest-error fonts
(Thuluth, Andalus). Two regimes emerge---on Akhbar no LLM strategy beats raw
OCR (clearest case of the over-correction threshold), whereas on Traditional
zero-shot achieves an 83.1\% relative reduction, the single largest improvement
in the study. The optimal strategy tracks baseline error rate, supporting
selective deployment.

\begin{table}[t]
\caption{Best LLM strategy per font (CER, diacritics stripped). $\Delta$ is
vs.\ OCR baseline in percentage points. The optimal strategy varies across
fonts, confirming that no single approach dominates universally.}
\label{tab:perfont}
\begin{center}
\small
\begin{tabular}{lrrrr}
\toprule
\textbf{Font} & \textbf{OCR} & \textbf{P2} & \textbf{Best} & \textbf{Best Strategy} \\
\midrule
Akhbar       &  2.12 &  2.71 &  2.46 & P6 conf\_self \\
Simplified   &  2.94 &  2.93 &  2.72 & P6 conf\_self \\
Traditional  &  3.63 &  2.96 &  \textbf{2.75} & P6 conf\_self \\
Arial        &  3.65 &  4.50 &  3.76 & P9 Error-Sig \\
Naskh        &  4.10 &  4.27 &  3.94 & P9 Error-Sig \\
Tahoma       &  5.82 &  6.13 &  5.76 & P9 Error-Sig \\
Thuluth      & 10.21 & 10.92 &  9.68 & P8 RAG \\
Andalus      & 12.12 & 12.27 & 10.41 & P8 RAG \\
\midrule
KHATT        & 34.24 & 33.25 & 33.13 & P9 Error-Sig \\
\bottomrule
\end{tabular}
\end{center}
\end{table}

\subsection{Experiment 2: Extended Prompt Design Study (13 Datasets)}

Experiment~2 crosses four prompt styles (base, conservative, error-pattern,
error-pattern+conservative) with two retrieval modes (zero-shot / BM25 RAG),
giving the eight trials in Table~\ref{tab:exp2_pats}. \textbf{T3 (conservative
zero-shot)} achieves the best PATS CER (5.27\%, $-$5.4\% relative vs.\ OCR),
beating all eight Experiment~1 phases---including P8 RAG (5.45\%)---without
retrieval overhead, showing that conservative framing is more effective than any
knowledge augmentation on clean typewritten text. \textbf{T2 (base+RAG)} is the
best overall configuration (5.44\% PATS, 33.82\% KHATT) and is selected for
Experiment~3. Error-pattern trials (T4--T8) are brittle; T7 fails catastrophically
(190.60\% CER) due to prompt-template tag injection.

\begin{table}[t]
\caption{Experiment~2 trial definitions and PATS-A01 macro-avg / KHATT results
(normal samples, diacritics stripped). $\dagger$ T7: template injection failure,
excluded. Best PATS in \textbf{bold}.}
\label{tab:exp2_pats}
\begin{center}
\small
\begin{tabular}{llccc}
\toprule
\textbf{\#} & \textbf{Style} & \textbf{RAG} & \textbf{PATS CER} & \textbf{KHATT CER} \\
\midrule
OCR & baseline  & --- & 5.57\% & 34.24\% \\
T1 & Base      & No  & 5.84\% & 33.25\% \\
T2 & Base      & Yes & 5.44\% & 33.82\% \\
\textbf{T3} & \textbf{Cons.} & \textbf{No} & \textbf{5.27\%} & 32.93\% \\
T4 & Cons.     & Yes & 27.02\% & 32.69\% \\
T5 & Err-pat.  & No  & 5.85\% & 32.90\% \\
T6 & Err-pat.  & Yes & 19.64\% & 32.64\% \\
T7$\dagger$ & Ep+Cons & No & 190.60\% & 51.64\% \\
T8 & Ep+Cons   & Yes & 46.56\% & 32.54\% \\
\bottomrule
\end{tabular}
\end{center}
\end{table}

For the four full-page datasets (Table~\ref{tab:exp2_fullpage}), T2 with runaway
correction (threshold $3.0\times$GT length) reduces the average CER from 86.50\%
to 68.18\% ($-$21.2\% absolute)---a far larger gain than on line-level PATS or
KHATT, driven mainly by the runaway corrector on domain-diverse page-level input.

\begin{table}[t]
\caption{Experiment~2 full-page dataset results: Qaari OCR baseline and
T2 (base+RAG) corrected CER$\dagger$ after runaway correction.
Qwen3-4B, normal samples, diacritics stripped.}
\label{tab:exp2_fullpage}
\begin{center}
\small
\begin{tabular}{lrr}
\toprule
\textbf{Dataset} & \textbf{OCR CER} & \textbf{T2 CER$\dagger$} \\
\midrule
KHATT-Paragraph  & 61.68\% & 45.06\% \\
Yarmouk-testing  & 49.85\% & 43.35\% \\
Muharaf-val      & 129.39\% & 89.78\% \\
Historical       & 105.10\% & 94.53\% \\
\midrule
\textbf{Full-page Avg} & \textbf{86.50\%} & \textbf{68.18\%} \\
\bottomrule
\end{tabular}
\end{center}
\end{table}

\subsection{Experiment 3: Multi-Model, Multi-OCR-Source, Multi-Domain}

Experiment~3 applies T2 across four LLMs, two OCR sources, and two held-out
benchmarks. In the model comparison (Table~\ref{tab:exp3_models}), Gemma-3-12B
is best on PATS (5.33\%) while Qwen3-4B is best on KHATT (33.82\%) and full-page
(68.18\%); Qwen3-14B suffers a KHATT runaway epidemic (175.69\% raw, recovered
to 35.55\%), so larger is not better within Qwen3.

\begin{table}[t]
\caption{Experiment~3A: four models on validation set with Qaari OCR
and T2. CER$\dagger$ = runaway-corrected. FP = full-page group avg.}
\label{tab:exp3_models}
\begin{center}
\small
\begin{tabular}{lcccr}
\toprule
\textbf{Model} & \textbf{Params} & \textbf{PATS$\dagger$} & \textbf{KHATT$\dagger$} & \textbf{FP$\dagger$} \\
\midrule
OCR Baseline & --- & 5.57\% & 34.24\% & 86.50\% \\
Qwen3-4B     & 4B  & 5.44\% & \textbf{33.82\%} & \textbf{68.18\%} \\
Qwen3-14B    & 14B & 5.60\% & 35.55\% & 88.51\% \\
Gemma-3-4B   & 4B  & 7.47\% & 36.14\% & 71.12\% \\
Gemma-3-12B  & 12B & \textbf{5.33\%} & 35.98\% & 70.01\% \\
\bottomrule
\end{tabular}
\end{center}
\end{table}

For the OCR-source comparison (Table~\ref{tab:exp3_ocr}), Gemma-3 VLM OCR is
excluded from PATS/KHATT (line-strip images are squashed by its preprocessor,
yielding hallucinations $>$170\% CER). On full-page inputs, Qaari and Gemma OCR
are comparable on average (86.50\% vs.\ 87.33\%) and converge after correction
(68.18\% vs.\ 64.28\%), so OCR-source choice matters less once the LLM corrector
is applied.

\begin{table}[t]
\caption{Experiment~3B: Qaari vs.\ Gemma-3 VLM OCR on full-page datasets
(Qwen3-4B, T2). CER$\dagger$ = runaway-corrected.
Gemma excluded for PATS/KHATT (line-strip preprocessing, invalid).}
\label{tab:exp3_ocr}
\begin{center}
\small
\begin{tabular}{lrrrr}
\toprule
\textbf{Dataset}
  & \multicolumn{2}{c}{\textbf{Qaari OCR}}
  & \multicolumn{2}{c}{\textbf{Gemma OCR}} \\
\cmidrule(lr){2-3}\cmidrule(lr){4-5}
& \textbf{OCR} & \textbf{Corr$\dagger$} & \textbf{OCR} & \textbf{Corr$\dagger$} \\
\midrule
KHATT-Paragraph  &  61.68\% & 45.06\% &  36.14\% & 35.82\% \\
Yarmouk-testing  &  49.85\% & 43.35\% &  95.76\% & 74.37\% \\
Muharaf-val      & 129.39\% & 89.78\% & 141.15\% & 72.85\% \\
Historical       & 105.10\% & 94.53\% &  76.26\% & 74.09\% \\
\midrule
\textbf{Avg} & \textbf{86.50\%} & \textbf{68.18\%} & \textbf{87.33\%} & \textbf{64.28\%} \\
\bottomrule
\end{tabular}
\end{center}
\end{table}

Correction generalises robustly to both held-out benchmarks
(Table~\ref{tab:exp3_gen}): on RDI-Test-Lines, Qaari 95.31\%~$\to$~85.11\%
($-$10.7\% abs) and Gemma 68.49\%~$\to$~61.13\%, with Gemma OCR outperforming
Qaari overall---its general VLM capability transfers better to unseen styles. On
Kitab, correction reduces CER by $-$17.3\% (Qaari) to $-$24.4\% (Gemma)
relative. The over-correction threshold reappears in new domains: low-CER
subsets (kitab-arabicocr 1.80\%, kitab-patsocr 4.79\%) are harmed, while
kitab-khattparagraph (205.27\%) gains most from runaway correction.

\begin{table}[t]
\caption{Experiment~3C: RDI-Test-Lines and Kitab benchmarks
(Qwen3-4B, T2, runaway-corrected). RDI Line Segmentation task excluded
(polygon-only GT, no CER computable). Kitab shows six representative
subsets plus overall. $\uparrow$ = over-correction (LLM correction harmful).}
\label{tab:exp3_gen}
\begin{center}
\small
\begin{tabular}{llrr}
\toprule
\textbf{Subset} & \textbf{Src} & \textbf{OCR CER} & \textbf{Corr$\dagger$} \\
\midrule
\multicolumn{4}{l}{\textit{RDI-Test-Lines (line recognition only)}} \\
LR-Handwritten  & Qaari &  98.82\% &  91.55\% \\
LR-Handwritten  & Gemma & \textbf{62.53\%} & \textbf{57.71\%} \\
LR-Manuscripts  & Qaari &  81.94\% &  79.00\% \\
LR-Manuscripts  & Gemma &  78.10\% &  70.36\% \\
LR-Typewritten  & Qaari & 105.19\% &  84.79\% \\
LR-Typewritten  & Gemma &  64.86\% &  55.31\% \\
\textbf{Overall}& Qaari &  95.31\% &  85.11\% \\
\textbf{Overall}& Gemma & \textbf{68.49\%} & \textbf{61.13\%} \\
\midrule
\multicolumn{4}{l}{\textit{Kitab Benchmark (selected subsets)}} \\
kitab-arabicocr      & Qaari &   1.80\% &   7.21\%$\uparrow$ \\
kitab-patsocr        & Qaari &   4.79\% &   5.98\%$\uparrow$ \\
kitab-hindawi        & Qaari &  31.46\% &  24.31\% \\
kitab-historyar      & Qaari &  63.48\% &  51.69\% \\
kitab-muharaf        & Qaari &  76.04\% &  58.72\% \\
kitab-khattparagraph & Qaari & 205.27\% &  90.60\% \\
\textbf{Kitab Avg}   & Qaari &  49.28\% & \textbf{40.73\%} \\
\textbf{Kitab Avg}   & Gemma &  79.45\% &  60.07\% \\
\bottomrule
\end{tabular}
\end{center}
\end{table}

% -------------------------------------------------------------------------
\section{Discussion}
% -------------------------------------------------------------------------

\subsection{The Over-Correction Threshold}

The font-dependent results point to a quantifiable \textit{over-correction
threshold}~$e^*$. Let~$r_f$ denote the model's fix rate and~$r_i$ its
introduction rate (fraction of correct tokens wrongly changed). The expected
output CER is $\text{CER}_{\text{out}} \approx e(1 - r_f) + (1-e)\,r_i$, so
correction is net beneficial when $e \cdot r_f > (1-e) \cdot r_i$, yielding
$e^* = r_i / (r_f + r_i)$. Empirically $r_i \approx 0.22$, $r_f \approx 0.76$ on
PATS-A01, giving $e^*\approx22\%$ in aggregate; the per-font crossing between
Akhbar (harmed) and Simplified (helped) places the effective threshold near
6--7\% CER, reflecting a drop in~$r_f$ at low error densities. Crucially, adding
context (confusion lists, diagnostics) signals to the model that the input is
defective and raises~$r_i$, pushing $e^*$ up and making correction harmful over
a wider range; conversely Phase~4 explicitly lowers~$r_i$, which is why it is the
only strategy to improve the lowest-error font and why \texttt{conf\_self}
synergises ($r_f$ up via P3, $r_i$ down via P4).

\subsection{Practical Implications and Limitations}

For clean text (CER $<$8\%), the conservative prompt outperforms every
knowledge-augmentation strategy; such output should be corrected conservatively
or left unchanged. Moderate-to-high-error output benefits from correction,
motivating selective deployment. RAG is effective only when the retrieval corpus
is domain-aligned; both this and the over-correction threshold generalise to
the held-out benchmarks. This study is limited to one primary LLM (Qwen3-4B)
and one OCR engine; CER/WER are proxies for downstream utility, which may rank
strategies differently, and self-reflective insights may not transfer across
models or domains.

% -------------------------------------------------------------------------
\section{Conclusion}
% -------------------------------------------------------------------------

Four findings emerge from our three-experiment study (thirteen datasets,
four models, two OCR sources, two held-out benchmarks; normal samples,
diacritics stripped).

\textbf{First, zero-shot LLM correction is net harmful on typewritten Arabic.}
Phase~2 raises PATS macro-average CER from 5.57\% to 5.84\% ($+$4.8\%), harming
seven of eight fonts. Traditional (3.63\%) is helped while Arial (3.65\%) is
harmed, showing error type and systematicity---not error rate alone---govern
the outcome.

\textbf{Second, domain-matched RAG is the only Experiment~1 strategy to beat
the OCR baseline on PATS.} Phase~8 (BM25) reaches 5.45\% ($-$0.12~pp vs.\ OCR)
and Phase~9 the best KHATT result (33.13\%). Ground-truth-anchored exemplars cut
false corrections on clean input; external-corpus RAG degrades severely,
confirming domain alignment as a prerequisite.

\textbf{Third, conservative prompting outperforms knowledge augmentation for
typewritten text.} Experiment~2's T3 (conservative zero-shot) reaches 5.27\%
PATS CER, better than all eight Experiment~1 phases including RAG, with no
retrieval overhead. Reducing intervention tendency beats supplying more
correction targets; error-pattern prompts (T4--T8) are brittle and T7 fails
catastrophically.

\textbf{Fourth, LLM correction generalises robustly to unseen domains.}
With T2, RDI-Test-Lines improves $-$10.7\% absolute and Kitab $-$17\% to
$-$24\% relative. The over-correction finding holds in new domains, so
multi-font, multi-domain evaluation is essential---single-dataset conclusions
are systematically misleading.

\begin{thebibliography}{00}

\bibitem{qaari}
Qaari OCR Engine, open-source Arabic OCR system.
Available: \url{https://github.com/quran/qaari}

\bibitem{arabic_ocr_survey}
A. Lorigo and V. Govindaraju, ``Offline Arabic handwriting recognition: a survey,''
\textit{IEEE Trans. Pattern Anal. Mach. Intell.}, vol. 28, no. 5, pp. 712--724, 2006.

\bibitem{pats}
M. Waly, A. Atwan, M. Cheriet, and O. Muda, ``PATS-A01: a large-scale Arabic
multi-font printed text corpus,'' in \textit{Proc. Int. Conf. Document Analysis
and Recognition (ICDAR)}, 2019.

\bibitem{khatt}
S. Mahmoud et al., ``KHATT: an open Arabic offline handwriting dataset,''
\textit{Pattern Recognition}, vol. 47, no. 3, pp. 1096--1112, 2014.

\bibitem{obeid2020}
O. Obeid et al., ``CAMeL tools: an open source python toolkit for
Arabic natural language processing,'' in \textit{Proc. LREC}, 2020, pp. 7022--7032.

\bibitem{qwen3}
Qwen Team, ``Qwen3 technical report,'' \textit{arXiv preprint arXiv:2505.09388},
2025.

\bibitem{afli2016}
H. Afli, L. Barrault, and H. Schwenk, ``OCR error correction using statistical
machine translation,'' \textit{Int. J. Comput. Linguistics Appl.}, vol. 7, no. 1,
pp. 175--191, 2016.

\bibitem{dong2018}
R. Dong, O. Lal, and G. Prud'hommeaux, ``Adapting sequence to sequence models
for text normalization in speech recognition,'' in \textit{Proc. INTERSPEECH}, 2018.

\bibitem{nguyen2021}
T. Nguyen, A. Jatowt, M. Coustaty, and A. Doucet, ``Survey of post-OCR processing
approaches,'' \textit{ACM Comput. Surv.}, vol. 54, no. 6, pp. 1--37, 2021.

\bibitem{llm_ocr_correction}
N. Kang et al., ``Large language models are better at OCR post-correction than you
might think,'' in \textit{Proc. COLING}, 2024.

\bibitem{arabic_bert_ocr}
W. Antoun, F. Baly, and H. Hajj, ``AraBERT: transformer-based model for Arabic language understanding,''
in \textit{Proc. LREC Workshop Arabic NLP}, 2020.

\bibitem{brown2020}
T. Brown \textit{et al.}, ``Language models are few-shot learners,''
\textit{Advances in Neural Information Processing Systems}, vol. 33, 2020.

\bibitem{wei2022}
J. Wei \textit{et al.}, ``Chain-of-thought prompting elicits reasoning in large
language models,'' \textit{Advances in Neural Information Processing Systems},
vol. 35, 2022.

\bibitem{lewis2020}
P. Lewis \textit{et al.}, ``Retrieval-augmented generation for knowledge-intensive
NLP tasks,'' \textit{Advances in Neural Information Processing Systems}, vol. 33,
2020.

\bibitem{shinn2023}
N. Shinn et al., ``Reflexion: language agents with verbal reinforcement learning,''
\textit{Advances in Neural Information Processing Systems}, vol. 36, 2023.

\bibitem{habash2010}
N. Habash, ``Introduction to Arabic natural language processing,''
\textit{Synthesis Lectures on HLT}, vol. 3, no. 1, pp. 1--187, 2010.

\bibitem{robertson2009bm25}
S. Robertson and H. Zaragoza, ``The probabilistic relevance framework: BM25 and beyond,''
\textit{Found. Trends Inf. Retr.}, vol. 3, no. 4, pp. 333--389, 2009.

\bibitem{romanov2017openiti}
M. Romanov and C. Sartori, ``OpenITI: a machine-readable corpus of Islamicate
texts,'' \textit{ADHO Digital Humanities}, 2019.

\end{thebibliography}

\end{document}
